\section{Применение формулы Литтла и соотношения в системах G/G/c}

Формула Литтла $L = \lambda W$ является универсальным инструментом, однако её применение требует четкого определения границ рассматриваемой системы («черного ящика»). Рассмотрим канонические примеры, демонстрирующие гибкость этого закона.

\begin{example}{Вариативность границ системы}
Рассмотрим стандартную одноканальную СМО. Закон Литтла можно применить к разным подсистемам, получая различные физические характеристики.

\begin{enumerate}
    \item \textbf{Система в целом (Очередь + Сервер).}
    Пусть «черный ящик» охватывает и буфер ожидания, и прибор обслуживания.
    
    \begin{center}
    \begin{tikzpicture}[scale=1.0, >=Stealth, thick]
        % Background box
        \draw[dashed, fill=Snow2, draw=DeepSkyBlue4] (-0.5, -0.7) rectangle (4.5, 1.2);
        \node[text=DeepSkyBlue4] at (2, 1.4) {Границы системы};

        % Queue
        \draw (0,0) -- (1.5,0) -- (1.5,0.8) -- (0,0.8);
        \foreach \x in {0.3, 0.6, 0.9, 1.2} \draw (\x,0) -- (\x,0.8);
        
        % Server
        \draw[fill=Snow2] (3,0.4) circle (0.4);
        \node at (3,0.4) {$\mu$};

        % Arrow connection
        \draw[->] (1.5, 0.4) -- (2.6, 0.4);

        % In/Out arrows
        \draw[->] (-1.5, 0.4) -- (0, 0.4) node[midway, above] {$\lambda$};
        \draw[->] (3.4, 0.4) -- (5, 0.4);
    \end{tikzpicture}
    \end{center}
    В этом случае $L$ — среднее число заявок в системе (в очереди и на обслуживании), $W$ — среднее время пребывания. Закон Литтла:
    $$ L = \lambda \cdot W. $$

    \item \textbf{Только очередь.}
    Рассмотрим подсистему, включающую только буфер ожидания. Сервер в данном контексте является внешней средой.
    
    \begin{center}
    \begin{tikzpicture}[scale=1.0, >=Stealth, thick]
        % Background box
        \draw[dashed, fill=Snow2, draw=DeepSkyBlue4] (-0.5, -0.5) rectangle (1.8, 1.2);
        \node[text=DeepSkyBlue4] at (0.65, 1.4) {Система};

        % Queue
        \draw (0,0) -- (1.5,0) -- (1.5,0.8) -- (0,0.8);
        \foreach \x in {0.3, 0.6, 0.9, 1.2} \draw (\x,0) -- (\x,0.8);
        
        % Server (outside)
        \draw[fill=Snow2, draw=gray] (3,0.4) circle (0.4);
        \node at (3,0.4) {$\mu$};
        
        % Arrow connection
        \draw[->] (1.5, 0.4) -- (2.6, 0.4);

        % In/Out arrows
        \draw[->] (-1.5, 0.4) -- (0, 0.4) node[midway, above] {$\lambda$};
        \draw[->] (3.4, 0.4) -- (5, 0.4);
    \end{tikzpicture}
    \end{center}
    Среднее число заявок в очереди $L_q$ связано со средним временем ожидания $W_q$ соотношением:
    $$ L_q = \lambda \cdot W_q. $$

    \item \textbf{Только сервер (канал обслуживания).}
    Рассмотрим систему, состоящую из одного прибора. Емкость такой системы равна 1. Пусть $p_0$ — вероятность того, что система пуста (сервер свободен). Тогда среднее число заявок в такой системе (обозначим $L_s$) равно вероятности застать сервер занятым:
    $$ L_s = 0 \cdot p_0 + 1 \cdot (1 - p_0) = 1 - p_0. $$
    
    \begin{center}
    \begin{tikzpicture}[scale=1.0, >=Stealth, thick]
        % Background box
        \draw[dashed, fill=Snow2, draw=DeepSkyBlue4] (2.2, -0.5) rectangle (3.8, 1.2);
        \node[text=DeepSkyBlue4] at (3, 1.4) {Система};

        % Queue (outside)
        \draw[draw=gray] (0,0) -- (1.5,0) -- (1.5,0.8) -- (0,0.8);
        
        % Server
        \draw[fill=Snow2] (3,0.4) circle (0.4);
        \node at (3,0.4) {$\mu$};
        
        % Arrow connection
        \draw[->] (1.5, 0.4) -- (2.6, 0.4) node[midway, above] {$\lambda$};
        
        % Out arrow
        \draw[->] (3.4, 0.4) -- (5, 0.4);
    \end{tikzpicture}
    \end{center}
    
    Среднее время пребывания в сервере равно среднему времени обслуживания: $W_s = \E[S] = \frac{1}{\mu}$.
    Применяя формулу Литтла $L_s = \lambda W_s$, получаем фундаментальную связь для одноканальной системы:
    $$ 1 - p_0 = \lambda \cdot \E[S] = \frac{\lambda}{\mu}. $$
    Это означает, что вероятность занятости сервера равна коэффициенту загрузки $\rho = \lambda / \mu$.
\end{enumerate}
\end{example}

\begin{note}
Для систем с ограниченной очередью и потерями формула Литтла применяется к потоку заявок, фактически попавших в систему. Если $p_B$ — вероятность отказа (блокировки), то эффективная интенсивность входного потока равна:
$$ \lambda_{\text{eff}} = \lambda (1 - p_B). $$
Тогда среднее число заявок в системе: $L = \lambda(1 - p_B) W$.
\end{note}

\subsection{Фундаментальные соотношения для систем G/G/c}

Рассмотрим систему общего вида $G/G/c$, где $G$ обозначает произвольное распределение интервалов поступления и времени обслуживания, а $c$ — число каналов. Введем базовые обозначения случайных величин и их моментов.

\begin{definition}{Параметры нагрузки и эффективность}
\begin{itemize}
    \item $\lambda$ — интенсивность входящего потока.
    \item $S$ — случайная величина времени обслуживания заявки.
    \item $\mu \doteq \frac{1}{\E[S]}$ — интенсивность обслуживания одного канала.
    \item $r \doteq \lambda \E[S] = \frac{\lambda}{\mu}$ — \textbf{интенсивность нагрузки} (измеряется в Эрлангах). Характеризует среднее число заявок, поступающих за время обслуживания одной заявки.
    \item $\rho \doteq \frac{r}{c} = \frac{\lambda}{c\mu}$ — \textbf{коэффициент загрузки} (utilization factor). Для стабильности системы необходимо $\rho < 1$.
\end{itemize}
\end{definition}

Определим характеристики пребывания заявки в системе:

\begin{itemize}
    \item $N, N_q$ — случайное число заявок в системе и в очереди соответственно.
    \item $L \doteq \E[N], \quad L_q \doteq \E[N_q]$ — средние значения (матожидания).
    \item $T, T_q$ — случайное время нахождения заявки в системе и в очереди.
    \item $W \doteq \E[T], \quad W_q \doteq \E[T_q]$ — средние времена.
\end{itemize}

Очевидно, что полное время пребывания складывается из времени ожидания и времени обслуживания:
$$ T = T_q + S. $$
В силу линейности математического ожидания, аналогичное равенство справедливо для средних величин:
$$ W = W_q + \E[S] = W_q + \frac{1}{\mu}. $$

\begin{statement}
Для любой системы $G/G/c$ в стационарном режиме среднее число заявок в системе $L$ связано со средним числом заявок в очереди $L_q$ и интенсивностью нагрузки $r$ соотношением:
$$ L = L_q + r. $$
\end{statement}

\begin{sproof}
Воспользуемся формулой Литтла для системы в целом ($L = \lambda W$) и полученным выше разложением для среднего времени пребывания:
$$
L = \lambda W = \lambda \left( W_q + \frac{1}{\mu} \right).
$$
Раскроем скобки:
$$
L = \underbrace{\lambda W_q}_{(*)} + \underbrace{\frac{\lambda}{\mu}}_{(**)}.
$$
Проанализируем слагаемые:
\begin{enumerate}
    \item По формуле Литтла для очереди, слагаемое $(*)$ равно среднему числу заявок в очереди: $\lambda W_q = L_q$.
    \item Слагаемое $(**)$ представляет собой интенсивность нагрузки: $\frac{\lambda}{\mu} = r$.
\end{enumerate}
Таким образом, приходим к искомому равенству $L = L_q + r$.
\end{sproof}

\begin{conseq}
Величину $r = L - L_q$ можно интерпретировать как среднее число заявок, находящихся \textit{на обслуживании} в данный момент времени.

В частном случае одноканальной системы ($c=1$) среднее число обслуживаемых заявок равно вероятности того, что сервер занят:
$$
r = 0 \cdot p_0 + 1 \cdot (1 - p_0) = 1 - p_0.
$$
Следовательно, для системы $G/G/1$:
$$ \rho = r = 1 - p_0. $$
Это подтверждает интуитивный факт: коэффициент загрузки одноканальной системы равен вероятности её занятости.
\end{conseq}

\section{Напоминание: Пуассоновский процесс и порядковые статистики}

Перед переходом к доказательству свойства PASTA необходимо зафиксировать определения базового стохастического процесса теории массового обслуживания и вспомогательный результат из математической статистики.

\begin{definition}{Пуассоновский процесс}
\textbf{Пуассоновской случайной величиной} $\xi \sim \text{Pois}(\lambda)$ называется величина с распределением:
$$ \P(\xi = n) = \frac{\lambda^n}{n!} e^{-\lambda}, \quad \E\xi = \mathbb{D}\xi = \lambda. $$

\textbf{Пуассоновским процессом} (потоком событий) $\{N(t), t \ge 0\}$ с интенсивностью $\lambda$ называется считающий процесс, удовлетворяющий следующим аксиомам:
\begin{enumerate}
    \item \textbf{Начальное условие:} $N(0) = 0$.
    \item \textbf{Ординарность:} Вероятность появления двух и более событий за малый промежуток времени пренебрежимо мала:
    $$ \P(N(t + \Delta t) - N(t) \ge 2) = o(\Delta t). $$
    Вероятность появления ровно одного события асимптотически линейна:
    $$ \P(N(t + \Delta t) - N(t) = 1) = \lambda \Delta t + o(\Delta t). $$
    \item \textbf{Отсутствие последействия:} Процесс имеет независимые приращения. Число событий на непересекающихся интервалах времени суть независимые случайные величины.
    \item \textbf{Стационарность (однородность):} Распределение числа событий зависит только от длины интервала, а не от его положения на оси времени.
\end{enumerate}
\end{definition}

Для дальнейшего анализа нам потребуется лемма о совместном распределении упорядоченных случайных величин.

\begin{lemma}{Совместная плотность порядковых статистик}
Пусть $X_1, \dots, X_n$ — независимые одинаково распределённые (i.i.d.) случайные величины с функцией плотности $f(x)$ и функцией распределения $F(x)$.
Рассмотрим вариационный ряд (упорядоченные по возрастанию значения):
$$ X_{(1)} \le X_{(2)} \le \dots \le X_{(n)}. $$
Тогда совместная плотность вероятности вектора $(X_{(1)}, \dots, X_{(n)})$ имеет вид:
$$
f_{X_{(1)}, \dots, X_{(n)}}(y_1, \dots, y_n) =
\begin{cases}
    n! \cdot \prod_{i=1}^n f(y_i), & \text{если } y_1 < y_2 < \dots < y_n, \\
    0, & \text{иначе.}
\end{cases}
$$
\end{lemma}

\begin{lproof}
Рассмотрим вероятность попадания упорядоченных величин в непересекающиеся интервалы. Пусть $y_1 < x_1 < y_2 < x_2 < \dots < y_n < x_n$. Обозначим искомое событие как $A$:
$$ A = \{y_1 < X_{(1)} \le x_1, \: y_2 < X_{(2)} \le x_2, \: \dots, \: y_n < X_{(n)} \le x_n\}. $$
Поскольку исходные величины $X_i$ непрерывны, вероятность совпадения значений равна нулю, и мы можем рассматривать строгое неравенство $X_{(1)} < \dots < X_{(n)}$.

Любая реализация, удовлетворяющая условию $A$, соответствует ровно одной перестановке исходных величин $X_1, \dots, X_n$, так как интервалы $(y_i, x_i]$ не пересекаются. Всего существует $n!$ возможных перестановок индексов $\pi \in \mathfrak{S}_n$.
Событие $A$ можно представить как объединение несовместных событий по всем перестановкам:
$$
A = \bigsqcup_{\pi \in \mathfrak{S}_n} \{ y_1 < X_{\pi(1)} \le x_1, \dots, y_n < X_{\pi(n)} \le x_n \}.
$$
В силу того, что $X_i$ независимы и одинаково распределены, вероятность каждого слагаемого в объединении одинакова и равна произведению вероятностей попадания в интервалы:
$$
\P(y_i < X_i \le x_i) = F(x_i) - F(y_i).
$$
Суммируя по всем $n!$ перестановкам, получаем:
$$
\P(A) = n! \cdot \prod_{i=1}^n (F(x_i) - F(y_i)).
$$
С другой стороны, эта же вероятность выражается через интеграл от совместной плотности:
$$
\P(A) = \int_{y_1}^{x_1} \dots \int_{y_n}^{x_n} f_{X_{(1)}, \dots, X_{(n)}}(u_1, \dots, u_n) \, du_n \dots du_1.
$$
Чтобы найти плотность, необходимо взять смешанную частную производную от вероятности по верхним пределам интегрирования $x_1, \dots, x_n$ (при стремлении $y_i \to x_i$):
$$
f_{X_{(1)}, \dots, X_{(n)}}(x_1, \dots, x_n) = \frac{\partial^n}{\partial x_1 \dots \partial x_n} \left[ n! \cdot \prod_{i=1}^n (F(x_i) - F(y_i)) \right].
$$
Поскольку каждый множитель $(F(x_i) - F(y_i))$ зависит только от своего $x_i$, производная произведения равна произведению производных. Учитывая, что $F'(x) = f(x)$, окончательно получаем:
$$
f_{X_{(1)}, \dots, X_{(n)}}(x_1, \dots, x_n) = n! \cdot \prod_{i=1}^n f(x_i),
$$
при условии упорядоченности аргументов.
\end{lproof}

\section{Связь пуассоновского потока и равномерного распределения}

Ключевым свойством пуассоновского процесса, лежащим в основе многих методов теории массового обслуживания, является его «абсолютная случайность» на любом конечном интервале. Математически это выражается в том, что при фиксированном числе событий их моменты распределены равномерно.

\begin{theorem}{Об условном распределении моментов наступления событий}
Пусть $\{N(t), t \ge 0\}$ — пуассоновский процесс. Зафиксируем интервал $[0, T]$ и предположим, что на нем произошло ровно $k$ событий ($N(T) = k$). Обозначим моменты этих событий через $0 < \tau_1 < \tau_2 < \dots < \tau_k < T$.

Тогда совместное распределение вектора $(\tau_1, \dots, \tau_k)$ совпадает с распределением $k$ \textbf{порядковых статистик} независимых случайных величин, равномерно распределенных на отрезке $[0, T]$.
\end{theorem}

\begin{tproof}
Рассмотрим совместную плотность распределения упорядоченных моментов времени $0 < t_1 < t_2 < \dots < t_k < T$. По определению плотности вероятности:
$$
f(t_1, \dots, t_k \mid N(T)=k) \, dt_1 \dots dt_k \approx \P(t_1 \le \tau_1 < t_1 + dt_1, \dots, t_k \le \tau_k < t_k + dt_k \mid N(T) = k).
$$
По формуле условной вероятности $\P(A \mid B) = \frac{\P(AB)}{\P(B)}$, вычислим числитель и знаменатель отдельно.

\textbf{1. Знаменатель (условие):}
Вероятность того, что на всем интервале $[0, T]$ произошло ровно $k$ событий, подчиняется распределению Пуассона:
$$ \P(N(T) = k) = \frac{(\lambda T)^k}{k!} e^{-\lambda T}. $$

\textbf{2. Числитель (совместное событие):}
Событие в числителе $AB$ означает, что траектория процесса реализовалась строго определенным образом. Разобьем интервал $[0, T]$ на последовательность непересекающихся промежутков:
\begin{itemize}
    \item На отрезке $[0, t_1)$ событий не было (0 событий).
    \item На малом отрезке $[t_1, t_1 + dt_1)$ произошло ровно 1 событие.
    \item На отрезке $[t_1 + dt_1, t_2)$ событий не было.
    \item $\dots$
    \item На отрезке $[t_k, t_k + dt_k)$ произошло ровно 1 событие.
    \item На оставшемся отрезке $[t_k + dt_k, T]$ событий не было.
\end{itemize}

В силу свойства \textbf{отсутствия последействия} (независимости приращений), вероятность такого сложного события равна произведению вероятностей на каждом участке.
Учитывая, что $\P(N(\Delta t)=0) = e^{-\lambda \Delta t}$ и $\P(N(dt)=1) \approx \lambda dt \cdot e^{-\lambda dt} \approx \lambda dt$, запишем:

$$
\P(AB) \approx \underbrace{e^{-\lambda t_1}}_{[0, t_1)} \cdot \underbrace{\lambda dt_1}_{[t_1, t_1+dt_1)} \cdot \underbrace{e^{-\lambda(t_2 - t_1)}}_{(t_1, t_2)} \cdot \dots \cdot \underbrace{\lambda dt_k}_{[t_k, t_k+dt_k)} \cdot \underbrace{e^{-\lambda(T - t_k)}}_{(t_k, T)}.
$$
(Здесь мы пренебрегли слагаемыми порядка $dt$ в показателях экспонент для интервалов пустоты).

Сгруппируем множители. Произведение $\lambda$ дает $\lambda^k dt_1 \dots dt_k$. Показатели экспонент складываются, образуя полную длину интервала:
$$
t_1 + (t_2 - t_1) + (t_3 - t_2) + \dots + (T - t_k) = T.
$$
Таким образом, вероятность совместного события равна:
$$ \P(AB) = \lambda^k e^{-\lambda T} \, dt_1 \dots dt_k. $$

\textbf{3. Итоговая плотность:}
Подставим найденные выражения в формулу условной вероятности:
$$
f(t_1, \dots, t_k \mid k) \, dt_1 \dots dt_k = \frac{\lambda^k e^{-\lambda T} \, dt_1 \dots dt_k}{\frac{(\lambda T)^k}{k!} e^{-\lambda T}}.
$$
Сокращая $\lambda^k$ и $e^{-\lambda T}$, получаем окончательный результат:
$$
f(t_1, \dots, t_k \mid k) = \frac{k!}{T^k}, \quad \text{при } 0 < t_1 < \dots < t_k < T.
$$
Это выражение в точности совпадает с плотностью $k$ порядковых статистик равномерного распределения $\mathcal{U}[0, T]$, полученной в предыдущей лемме.
\end{tproof}

\subsection{Интерпретация}

Полученный результат позволяет дать наглядную интерпретацию поведения пуассоновского потока. Представим, что мы бросаем $k$ дротиков в временную ось $[0, T]$ абсолютно независимо и «вслепую» (равномерное распределение).

\begin{itemize}
    \item Пусть координаты попадания дротиков: $U_1, U_2, \dots, U_k$. Это независимые, но \textbf{неупорядоченные} величины ($U_1$ может быть больше $U_2$).
    \item Реальные события в потоке всегда происходят последовательно. Согласно доказанной теореме, моменты событий пуассоновского потока $\tau_1, \dots, \tau_k$ соответствуют \textbf{отсортированным} координатам дротиков $U_{(1)} \le U_{(2)} \le \dots \le U_{(k)}$.
\end{itemize}

\begin{center}
\begin{tikzpicture}[scale=1.5]
    % Axis
    \draw[->, thick] (0,0) -- (6,0) node[right] {$t$};
    \node at (0,-0.3) {$0$};
    \node at (5,-0.3) {$T$};
    \draw (5, 0.1) -- (5, -0.1);

    % Random points (Unordered)
    \foreach \x/\l in {1.2/U_2, 3.5/U_1, 4.2/U_3} {
        \draw[DeepSkyBlue4, thick] (\x, 0.1) -- (\x, -0.1);
        \node[DeepSkyBlue4, above] at (\x, 0.2) {$\l$};
    }
    
    % Ordered points
    \node at (2.5, -0.8) {$\Updownarrow$ Сортировка};
    
    % Axis 2
    \draw[->, thick] (0,-1.5) -- (6,-1.5) node[right] {$t$};
    \node at (0,-1.8) {$0$};
    \node at (5,-1.8) {$T$};
    \draw (5, -1.4) -- (5, -1.6);
    
    \foreach \x/\l in {1.2/\tau_1, 3.5/\tau_2, 4.2/\tau_3} {
        \draw[Red3, thick] (\x, -1.4) -- (\x, -1.6);
        \node[Red3, below] at (\x, -1.6) {$\l$};
    }
\end{tikzpicture}
\end{center}

\textbf{Вывод:} Положение событий пуассоновского потока на оси времени «максимально случайно». У событий нет склонности к группированию или, наоборот, к регулярному следованию друг за другом. Это обеспечивает «непредвзятость» наблюдения системы в моменты поступления заявок.

\section{Свойство PASTA}

Доказанная теорема является фундаментом для одного из самых известных эвристических принципов теории массового обслуживания — свойства PASTA (\textit{Poisson Arrivals See Time Averages}).

\begin{theorem}{PASTA}
Пусть $X(t)$ — случайный процесс, описывающий состояние системы (например, число заявок в очереди). Если входящий поток заявок является пуассоновским, то предельное распределение состояния системы, которое «видит» поступающая заявка, совпадает с предельным распределением состояния системы во времени:
$$
\lim_{n \to \infty} \P(X(\tau_n - 0) = j) = \lim_{t \to \infty} \P(X(t) = j).
$$
\textit{«Пуассоновские заявки видят временн\'{ы}е средние (видят среднее поведение системы)».}
\end{theorem}

\textbf{Физический смысл:}
Поскольку пуассоновские события разбросаны по оси времени равномерно (как показано в теореме 6), они сканируют состояние системы в случайные моменты времени. Сам факт связи с равномерным распределением означает, что пуассоновский поток не имеет "предпочтительных" моментов времени или внутренней периодичности. Это означает, что входящий поток не «подгадывает» моменты, когда система пуста или перегружена (как это мог бы делать регулярный поток, синхронизированный с обслуживанием). Поэтому, в частности, средние характеристики системы, наблюдаемые в моменты поступления заявок, совпадают со средними характеристиками системы в целом.

\begin{note}
Это свойство не выполняется для других потоков. Например, если заявки поступают строго периодически (детерминированный поток $D$) каждые 10 минут, а обслуживание занимает ровно 9 минут, то каждая новая заявка всегда будет видеть систему свободной ($X(\tau_n) = 0$), хотя 90\% времени система занята. Пуассоновский поток лишен такой «синхронизации».
\end{note}

\section{Система M|M|1}

Рассмотрим классическую одноканальную систему массового обслуживания с простейшим входным потоком и показательным временем обслуживания. В нотации Кендалла она обозначается как $M|M|1|\infty$.

\begin{statement}
Пусть $X(t)$ — число заявок в системе в момент времени $t$. Поскольку входной поток пуассоновский, а обслуживание экспоненциальное, процесс $X(t)$ является \textbf{процессом рождения и гибели} со следующими параметрами:
\begin{itemize}
    \item Интенсивность рождения (поступления): $\lambda_n = \lambda$ для всех $n \ge 0$.
    \item Интенсивность гибели (обслуживания): $\mu_n = \mu$ для всех $n \ge 1$ (при $n=0$ обслуживание невозможно).
\end{itemize}
\end{statement}

\begin{note}
Важно отметить, что свойства конечных систем (например, $M|M|1|K$) нельзя автоматически переносить на бесконечные системы путем предельного перехода $K \to \infty$ без дополнительного анализа устойчивости. В конечной цепи стационарное распределение существует при любых $\lambda$ и $\mu$. В бесконечной цепи — только при условии $\lambda < \mu$, как будет показано далее.
\end{note}

\subsection*{Граф состояний и уравнения баланса}

Построим диаграмму интенсивностей переходов. Состояния системы соответствуют числу заявок $S = \{0, 1, 2, \dots\}$.

\begin{center}
    \begin{tikzpicture}[
        scale=1.5,
        >={Stealth[length=3mm]},
        state/.style={circle, draw, minimum size=0.8cm, thick},
        every node/.style={font=\small}
    ]
        % Состояния
        \node[state] (s0) at (0,0) {0};
        \node[state] (s1) at (2,0) {1};
        \node[state] (s2) at (4,0) {2};
        \node (dots) at (6,0) {\Large $\dots$};
        \node[state] (sn) at (8,0) {$n$};
        \node (dots2) at (10,0) {\Large $\dots$};

        % Переходы вперед (лямбда)
        \draw[->, thick] (s0) to [bend left=40] node[midway, above] {$\lambda$} (s1);
        \draw[->, thick] (s1) to [bend left=40] node[midway, above] {$\lambda$} (s2);
        \draw[->, thick] (s2) to [bend left=40] node[midway, above] {$\lambda$} (dots);
        \draw[->, thick] (dots) to [bend left=40] node[midway, above] {$\lambda$} (sn);
        \draw[->, thick] (sn) to [bend left=40] node[midway, above] {$\lambda$} (dots2);

        % Переходы назад (мю)
        \draw[->, thick] (s1) to [bend left=40] node[midway, below] {$\mu$} (s0);
        \draw[->, thick] (s2) to [bend left=40] node[midway, below] {$\mu$} (s1);
        \draw[->, thick] (dots) to [bend left=40] node[midway, below] {$\mu$} (s2);
        \draw[->, thick] (sn) to [bend left=40] node[midway, below] {$\mu$} (dots);
        \draw[->, thick] (dots2) to [bend left=40] node[midway, below] {$\mu$} (sn);
    \end{tikzpicture}
\end{center}

Для нахождения стационарных вероятностей $p_n = \lim_{t \to \infty} \P(X(t)=n)$ воспользуемся принципом локального баланса («поток через разрез между состояниями должен быть равен нулю»).

Для любого разреза между состояниями $n$ и $n+1$ поток вероятности слева направо (рождение) должен уравновешиваться потоком справа налево (гибель):
$$
\lambda p_n = \mu p_{n+1}, \quad n \ge 0.
$$
Отсюда получаем рекуррентное соотношение:
$$
p_{n+1} = \frac{\lambda}{\mu} p_n.
$$
Введем коэффициент загрузки системы $\rho \doteq \frac{\lambda}{\mu}$. Тогда:
$$
p_{n+1} = \rho p_n \implies p_n = \rho^n p_0.
$$

\subsection*{Условие эргодичности и финальное распределение}

Для существования стационарного распределения необходимо выполнение условия нормировки $\sum p_n = 1$. Подставим выражение через $p_0$:
$$
\sum_{n=0}^{\infty} p_n = \sum_{n=0}^{\infty} p_0 \rho^n = p_0 \sum_{n=0}^{\infty} \rho^n = 1.
$$
Данный ряд является геометрической прогрессией. Он сходится тогда и только тогда, когда знаменатель прогрессии по модулю меньше единицы:
$$
\rho < 1 \iff \lambda < \mu.
$$
Если условие выполнено, сумма ряда равна $\frac{1}{1-\rho}$, откуда находим вероятность простоя системы:
$$
p_0 \cdot \frac{1}{1-\rho} = 1 \implies p_0 = 1 - \rho.
$$
Заметим, что этот результат совпадает с формулой $p_0 = 1 - \rho$, выведенной ранее для системы общего вида $G|G|1$ с использованием закона Литтла.

\begin{theorem}{Распределение числа заявок в M|M|1}
Если интенсивность поступления заявок меньше интенсивности обслуживания ($\lambda < \mu$), то число заявок в системе $M|M|1$ в стационарном режиме имеет \textbf{геометрическое распределение} с параметром $1-\rho$:
$$
p_n = (1 - \rho)\rho^n, \quad n = 0, 1, 2, \dots
$$
\end{theorem}

\begin{conseq}
Если $\rho \ge 1$, очередь растет неограниченно, стационарного режима не существует, и $p_n \to 0$ для любого фиксированного $n$ (вероятность застать конечное число заявок стремится к нулю).
\end{conseq}